\section{Scaling Law 与 Emergence：平滑总损失如何隐藏任务突变}

前一章把训练目标拆成潜在任务混合；本章研究规模轴。Jason 的第二、第三个直觉并不矛盾：增加计算量时，总体 cross-entropy 可以稳定下降，而某个 benchmark 的 accuracy 仍可能长期为零后突然上升。要理解这种现象，必须同时看连续概率、任务权重和离散评分规则。

\subsection{Scaling 的最小操作定义}

本节先说明横轴。课堂把训练 compute 粗略理解为模型规模与数据量的乘积；真实训练还取决于序列长度、前后向计算、硬件利用率和训练配方，但这个简化足以表达“投入更多有效计算”的方向。重要的是比较同一训练范式下的一系列点，而不是孤立比较两个模型名称。

\lecturefigure{jason-slide-008.jpg}{直觉二：扩大模型与数据所代表的训练计算会可靠降低 loss}{Jason Wei official deck, p.~8}

一个常见近似是
\[
C\propto N D,
\]
其中 $C$ 表示训练计算量，$N$ 表示非 embedding 参数规模，$D$ 表示训练 token 数；比例常数吸收每个 token 的前向与反向运算。这个式子不是完整成本模型，但能提醒研究者：只增参数而不增数据，或只增数据而模型容量不足，都可能偏离 compute-optimal 区域。

\lecturefigure{jason-slide-009.jpg}{语言模型 loss 随计算量呈可预测的幂律趋势}{Jason Wei official deck, p.~9}

典型 scaling-law 拟合写成
\[
L(C)=L_{\infty}+A C^{-\alpha},\qquad A>0,\ \alpha>0.
\]
其中 $L(C)$ 是给定 compute 下的验证损失，$L_{\infty}$ 是不可约或暂未消除的损失下界，$A$ 控制曲线高度，$\alpha$ 控制随 compute 改善的斜率。对数坐标下，幂律区间接近直线，因此可以用较小实验估计更大训练的预期 loss。

\begin{knowledgebox}{读图：先确认坐标、范围与是否饱和}
横轴跨越多个数量级，纵轴是越低越好的 loss。最关键的视觉证据不是某个点领先，而是趋势在很大范围内没有明显向水平线弯折。它支持“在观察区间内继续扩大计算仍有收益”，却不证明趋势永远持续，也不说明收益足以抵消训练成本。
\end{knowledgebox}

\teachervoice{Jason 在 00:09:36--00:10:44 提醒，这条关系被称为 ``law''，是因为它跨越约七个数量级仍保持可预测，而不是因为它像物理定律一样永不失效。研究报告应同时给拟合区间和外推不确定性。}

\subsection{为什么 scaling 有效：证据与猜测要分开}

上一节是经验关系，本节讨论可能机制。课堂给出的两种解释都被明确称为 hand-wavy：更大的模型能记住更多长尾事实，也能在一次 forward pass 中表示更复杂的函数。它们是有用假设，但不能仅凭 loss 曲线得到确认。

\lecturefigure{jason-slide-010.jpg}{小模型选择性记忆，大模型容纳长尾并执行更复杂计算的直觉}{Jason Wei official deck, p.~10}

\begin{warningbox}{读图时容易犯的因果错误}
“更大模型 loss 更低”不能单独证明“下降全部来自 memorization”，也不能证明“深度负责推理、宽度负责事实”。参数量、数据覆盖、优化稳定性和训练时间同步变化。要检验机制，需要控制变量、构造反事实数据，或直接测量不同频率事实和不同计算深度任务的变化。
\end{warningbox}

\teachervoice{在 00:11:13--00:12:55，Jason 两次说这些只是 hand-wavy answers。保留这层不确定性比给出一个听起来完整的机制更重要：课程事实是 scaling 关系稳健；“长尾记忆 + 更复杂函数”是解释候选。}

\subsection{总体 loss 是任务损失的加权混合}

现在回到前一章的 massive multi-task 视角。若语料中不同 token 分别需要语法、事实、情感或数学，那么总体 loss 就像这些子分布损失的加权平均。容易任务可以接近饱和，难任务仍在快速改善；因此总体曲线的平滑并不会强迫每个组成部分都平滑。

\lecturefigure{jason-slide-011.jpg}{直觉三：总体 loss 平滑下降，单个下游任务仍可能呈 emergent pattern}{Jason Wei official deck, p.~11}

\lecturefigure{jason-slide-012.jpg}{把总体 loss 分解成 grammar、world knowledge、math 等任务损失}{Jason Wei official deck, p.~12}

概念性地写，
\[
L_{\mathrm{all}}(C)=\sum_{k=1}^{K} w_k L_k(C),
\qquad w_k\ge 0,\quad \sum_{k=1}^{K}w_k=1.
\]
其中 $K$ 是潜在任务族数量，$w_k$ 是该任务在评估 token 中的权重，$L_k(C)$ 是 compute 为 $C$ 时的任务损失。即使某个困难任务 $L_j$ 在一个窄区间快速下降，只要 $w_j$ 很小，总体曲线仍可能看起来平滑。

\begin{knowledgebox}{什么是 emergent ability}
本讲采用行为层定义：小模型在某任务上接近随机或零分，而较大模型达到明显更好成绩。这个定义描述\textbf{评测曲线形状}，不自动说明模型内部突然产生了全新模块。若评分函数有阈值、exact match 或多步全对要求，连续概率改善也会被映射成表面突变。
\end{knowledgebox}

\subsection{容易任务饱和，困难任务晚启动}

上一小节给出混合公式，本节用图读出组成曲线。语法在较小规模可能已足够好，继续扩大模型只带来微小改进；数学等困难任务则可能在更大规模才跨过可用区间。读图时要比较斜率和饱和位置，而不是把每条线都理解成独立模块。

\lecturefigure{jason-slide-013.jpg}{总体 loss、已饱和任务与困难任务的不同 scaling 形状}{Jason Wei official deck, p.~13}

若任务得分由连续置信度 $q(C)$ 经阈值 $\tau$ 转换，
\[
S(C)=\mathbf{1}\{q(C)\ge \tau\},
\]
其中 $S(C)$ 是离散得分，$q(C)$ 是随 compute 平滑变化的成功置信度，$\tau$ 是评分阈值。即使 $q(C)$ 平滑上升，$S(C)$ 也会在跨越 $\tau$ 时从 0 跳到 1。这是 emergence 的一种来源，但不是唯一来源。

\lecturefigure{jason-slide-014.jpg}{BIG-Bench 任务包含 smooth、flat、inverse、noisy 与 emergent-looking 曲线}{Jason Wei official deck, p.~14}

\begin{knowledgebox}{读图：任务分布比“是否涌现”二分法更有信息}
图中约 29\% 任务平滑改善，22\% 在观测范围内近乎平坦，2\% 呈 inverse scaling，13\% 与规模关系不清，33\% 呈 emergent-looking。百分比依赖模型集合、任务和 metric；最稳健的结论是任务曲线具有多种形状，不能用总体 loss 的单一趋势替代逐任务测量。
\end{knowledgebox}

\teachervoice{Jason 在 00:16:54--00:19:24 强调：若只训练到突变点之前，研究者会错误预测该能力永远不会出现。Q\&A 又补充，缺少密集的中间模型点使“突变前是否已有信号”很难判断。}

\subsection{Prompting 中的 emergence 与 metric caveat}

任务曲线之外，prompt 本身也能暴露尺度差异。小模型可能无法从几个演示中抽取映射规则，大模型则能利用上下文完成新输入；但观察到 exact-match 跃迁后，仍应检查 token-level probability、连续 metric 与 prompt 稳健性，避免把评分函数的几何形状误写成内部机制。

\lecturefigure{jason-slide-015.jpg}{Prompt 中的输入输出示例可触发规模相关的 in-context behavior}{Jason Wei official deck, p.~15}

\begin{warningbox}{“Emergence 是 mirage”与“能力是假的”不是同一句话}
Q\&A 中 Jason 认可：改变 metric 会让曲线形状不同；但他仍认为大模型展示的有用行为是真实的。严谨表述应分三层：模型输出分布是否连续变化、评测分数是否突然变化、用户可完成的任务是否出现质变。三层可以同时成立而不矛盾。
\end{warningbox}

\teachervoice{00:29:53--00:30:21 的回答没有给出简单阵营标签，而是建议读论文并自己判断。讲义因此不把 emergence 处理成信仰问题，而把 metric sensitivity 当作实验设计的一部分。}

\subsection{本章小结}

Scaling law 描述总体损失在观测区间内随计算量稳定改善；emergent-looking curve 描述特定任务经过 metric 映射后的行为。两者可由任务混合、难度差异、阈值评分和稀疏采样共同产生。研究者应同时报告连续损失、任务分数、模型规模点与误差条，并把机制解释标明为证据、假设还是个人直觉。

